\begin{abstract}
    La compresión fractal representa una imagen mediante un sistema de funciones iteradas (IFS) basado en la autosimilitud entre bloques rango y dominio. Su rendimiento depende del contenido y la codificación es mucho más costosa que la decodificación.En este trabajo se evalúa la influencia del tipo de imagen y de su autosimilitud percibida sobre la calidad (PSNR, MSE, MAE), la tasa (bpp y relación de compresión) y los tiempos de codificación/decodificación, comparando tres geometrías de bloques con JPEG y PNG. La autosimilitud permite reconstruir mejor las imagenes y el contenido tipográfico es el caso más desfavorable. La tasa queda dominada por la geometría y el formato textual del codebook. JPEG supera al método fractal en tasa-distorsión y tiempo de codificación. PNG refleja la redundancia local, no necesariamente la autosemejanza geométrica.
    \keywords{compresión fractal, IFS, autosimilitud, JPEG, PNG}
    \end{abstract}
